\begin{center}
\hypertarget{native-061-article044}{}\pdfbookmark[2]{44.}{article-044}44.
\end{center}
Motum \emph{hyperbolicum} eo brevius absolvere licebit, quoniam methodo ei quam hactenus pro motu elliptico exposuimus prorsus analoga tractandus est. Aequationem inter tempus $t$ atque quantitatem auxiliarem $u$ forma sequente exhibemus:
\[
\begin{aligned}
&(e-1)\left\{\tfrac{1}{20}\left(u-\frac{1}{u}\right)+\tfrac{9}{10}\log u\right\}\\
&\quad+\left(\tfrac{1}{10}+\tfrac{9}{10}e\right)\left\{\tfrac{1}{2}\left(u-\frac{1}{u}\right)-\log u\right\}=kt\left(\frac{e-1}{q}\right)^{\frac{3}{2}},
\end{aligned}
\]
\noindent ubi logarithmi sunt hyperbolici, atque $\tfrac{1}{20}\left(u-\dfrac{1}{u}\right)+\tfrac{9}{10}\log u$ quantitas ordinis primi, $\tfrac{1}{2}\left(u-\dfrac{1}{u}\right)-\log u$ quantitas ordinis tertii, simulac $\log u$ tamquam quantitas parva ordinis primi spectatur. Statuendo itaque
\[
\begin{aligned}
\frac{6\left\{\tfrac{1}{2}\left(u-\frac{1}{u}\right)-\log u\right\}}{\tfrac{1}{20}\left(u-\frac{1}{u}\right)+\tfrac{9}{10}\log u}&=4A,\\[.7em]
\frac{\tfrac{1}{20}\left(u-\frac{1}{u}\right)+\tfrac{9}{10}\log u}{2\BareRoot{A}}&=B,
\end{aligned}
\]
\noindent\SourceLeaf{062}erit $A$ quantitas ordinis secundi, $B$ autem ab unitate differentia ordinis quarti discrepabit. Aequatio nostra tunc formam sequentem induet:
\[
\text{2)}\qquad B\left\{2(e-1)A^{\frac{1}{2}}+\tfrac{2}{15}(1+9e)A^{\frac{3}{2}}\right\}=kt\left(\frac{e-1}{q}\right)^{\frac{3}{2}},
\]
\noindent quae aequationi 1) art. 37 prorsus analoga est. Statuendo porro $\left(\dfrac{u-1}{u+1}\right)^2=T$, erit $T$ ordinis secundi, et per methodum serierum infinitarum invenietur
\[
\begin{aligned}
\frac{A}{T}&=1+\tfrac{4}{5}A+\tfrac{8}{175}A^2-\tfrac{8}{525}A^3\\
&\quad+\tfrac{1896}{336875}A^4-\tfrac{28744}{13138125}A^5+\text{etc.}
\end{aligned}
\]
\noindent Quamobrem ponendo $\dfrac{A}{T}=1+\tfrac{4}{5}A+C$, erit $C$ quantitas ordinis quarti, atque $A=\dfrac{(1+C)T}{1-\tfrac{4}{5}T}\mathord{\raisebox{.35ex}{.}}$ Denique pro radio vectore ex aequ. VII art. 21 facile sequitur
\[
r=\frac{q}{(1-T)\cos\tfrac{1}{2}v^2}=\frac{(1+\tfrac{4}{5}A+C)q}{(1-\tfrac{1}{5}A+C)\cos\tfrac{1}{2}v^2}\mathord{\raisebox{.35ex}{.}}
\]
